% Bài báo PBL3 nhóm 8 – bản nháp. Mẫu IEEE conference, tiếng Việt.
% Biên dịch: tectonic -X compile main.tex   (hoặc Overleaf, Compiler = XeLaTeX)
% \canlam{...} (chữ đỏ) = chỗ còn thiếu số liệu / cần kiểm tra; bỏ hết trước khi nộp.
% Số liệu mô phỏng lấy từ laptop/README.md (mục Cảnh báo), firmware/thu_mo_hinh/README.md.
\documentclass[conference]{IEEEtran}
\usepackage{fontspec}
% TeX Gyre: giống Times/Helvetica/Courier của IEEE, đủ dấu tiếng Việt; gọi theo tên file để chạy được cả trên
% tectonic lẫn Overleaf
\setmainfont{texgyretermes}[Extension=.otf,UprightFont=*-regular,BoldFont=*-bold,ItalicFont=*-italic,BoldItalicFont=*-bolditalic]
\setsansfont{texgyreheros}[Extension=.otf,UprightFont=*-regular,BoldFont=*-bold,ItalicFont=*-italic,BoldItalicFont=*-bolditalic]
\setmonofont{texgyrecursor}[Extension=.otf,UprightFont=*-regular,BoldFont=*-bold,ItalicFont=*-italic,BoldItalicFont=*-bolditalic]
\usepackage{amsmath,booktabs,array,xcolor,graphicx,url}
\usepackage{tikz}
\usetikzlibrary{positioning,arrows.meta,calc}
\newcommand{\canlam}[1]{\textcolor{red}{[#1]}}
% tên phần theo tiếng Việt
\renewcommand{\abstractname}{Tóm tắt}
\renewcommand{\IEEEkeywordsname}{Từ khoá}
\renewcommand{\tablename}{Bảng}
\renewcommand{\figurename}{Hình}
\renewcommand{\refname}{Tài liệu tham khảo}

\begin{document}

\title{Giám sát an toàn người ở nhà một mình\\bằng thông tin trạng thái kênh Wi‑Fi trên ESP32}

\author{\IEEEauthorblockN{Lê Minh Chiến, Nguyễn Thị Hà Tiên, Nguyễn Quang Phúc}
\IEEEauthorblockA{Khoa Điện tử và Trí tuệ nhân tạo, Trường Đại học Bách khoa – Đại học Đà Nẵng\\
Lớp 23KTMT2 · Nhóm 8 PBL3 Mạng máy tính}}

\maketitle

\begin{abstract}
Người sống một mình có thể bị ngã, bất động lâu hoặc thở bất thường mà không ai biết để giúp kịp. Bài báo đề xuất một
hệ thống giám sát không dùng camera, không cần đeo thiết bị, chỉ dựa vào thông tin trạng thái kênh (CSI) Wi‑Fi đo bằng
các board ESP32 giá rẻ: một ESP32‑C3 phát gói ESP‑NOW và hai ESP32‑S3 thu CSI. Xử lý chia thành bốn tầng: trạng thái
phòng, số người, động tác (gồm ngã) và nhịp thở; kết quả đưa vào bộ cảnh báo theo sự kiện có bước xác nhận. Mỗi
ESP32‑S3 tự chạy mô hình XGBoost trên bộ đặc trưng gọn (19 và 25 cột) rồi gộp kết quả hai board theo cách gộp muộn.
Trên 40 buổi mô phỏng (4,3 giờ), hệ thống phát hiện 44/53 lần ngã (83\,\%) với 5,1 lần báo ngã giả mỗi giờ; mô hình
động tác chạy trên ESP32‑S3 thật mất khoảng 82\,ms mỗi window. \canlam{Bổ sung kết quả so với baseline và kết quả trên
dữ liệu thật.}
\end{abstract}

\begin{IEEEkeywords}
Wi‑Fi sensing, CSI, ESP32, phát hiện ngã, nhịp thở, XGBoost, edge AI
\end{IEEEkeywords}

% =====================================================================================================================
\section{Giới thiệu}
Số người cao tuổi sống một mình ngày càng tăng \canlam{dẫn số liệu}. Với họ, các sự cố như ngã mà không tự đứng dậy
được, ngất, đột quỵ hay rối loạn hô hấp đặc biệt nguy hiểm vì không có ai bên cạnh. Camera giải quyết được bài toán nhưng
xâm phạm riêng tư; vòng tay, đồng hồ thông minh thì phụ thuộc vào việc người dùng nhớ đeo và sạc.

Wi‑Fi sensing tận dụng việc cơ thể người làm thay đổi đường lan truyền của sóng giữa máy phát và máy thu. Thông tin
trạng thái kênh (Channel State Information – CSI) mô tả biên độ và pha trên từng sóng mang con OFDM, nhạy với cả cử động
lớn (đi lại, ngã) lẫn dao động nhỏ của lồng ngực khi thở. Các board ESP32 cho phép đọc CSI với chi phí rất thấp.

Đóng góp của bài báo gồm:
\begin{itemize}
  \item Một pipeline bốn tầng biến CSI thành \emph{sự kiện cảnh báo} (ngã, ngã khẩn cấp, bất động lâu, thở bất thường,
    ngưng thở), với bước xác nhận 10\,s và cơ chế xử lý khi kênh bị thiết bị khác chiếm.
  \item Bộ đặc trưng gọn (19 cột cho window 6\,s, 25 cột cho window 3\,s) thay cho thống kê trên từng sóng mang con,
    đủ nhỏ để mỗi ESP32‑S3 tự chạy mô hình – bước mà \cite{natarajan2023} mới đề xuất – và kiểm chứng trên board
    thật.
  \item Gộp muộn kết quả của hai board theo nguyên tắc an toàn: ngã chỉ cần một board thấy, chế độ nhiều người cần cả
    hai cùng thấy.
  \item Đánh giá theo sự kiện (tỉ lệ bắt được, báo giả mỗi giờ, độ trễ) và so sánh với các phương pháp gốc.
    \canlam{hoàn thiện khi có bảng baseline}
\end{itemize}

% =====================================================================================================================
\section{Các nghiên cứu liên quan}
\textbf{Phát hiện chuyển động bằng RSSI và CSI trên thiết bị IoT.} Natarajan và cộng sự \cite{natarajan2023} thu RSSI
và CSI bằng laptop, điện thoại và ESP32 (nghe khung của router) tại hai toà nhà, tạo sáu bộ dữ liệu; mỗi window không
chồng lấn được mô tả bằng bốn đại lượng độ phân tán (phương sai, MAD, range, IQR) trên từng sóng mang con, sau khi khử
nhiễu wavelet Daubechies. Mô hình hai tầng gồm Random Forest, Gradient Boosting, XGBoost ở tầng 1 và bộ bỏ phiếu có
trọng số (cứng và mềm) ở tầng 2. Độ chính xác kiểm định chéo trung bình đạt tới 99,7\,\% khi có tầm nhìn thẳng và
97,8\,\% khi xuyên tường; RSSI xuyên tường có bộ chỉ đạt 75\,\%, và giảm ít nhất 21,2\,\% khi có nhiều người, trong khi
CSI giữ ổn định. Các mô hình chiếm dưới 1\,MB và được đề xuất là có thể đưa lên thiết bị biên, nhưng bài chưa chạy mô
hình trên ESP32 và chỉ phân biệt có người đi lại hay phòng trống.

\textbf{Dữ liệu công khai trên ESP32‑S3.} WiFall \cite{wifall} do router nhà phát và ESP32‑S3 thu (100\,Hz, 52 sóng mang
con), gồm đi, nhảy, ngồi, đứng dậy và năm kiểu ngã của 10 tình nguyện viên. Trong kịch bản khác người (train trên 9
người, một mẫu mỗi lớp của người thứ 10), KNN‑MMD đạt 77,62\,\% phát hiện ngã và 75,30\,\% nhận dạng động tác; ResNet18
train cùng miền đạt 91,88\,\% và 70,50\,\%. CSI‑BERT2 \cite{csibert2} báo 94,79\,\% phát hiện ngã, 88,59\,\% nhận dạng
động tác trên WiFall và 94,32\,\% đếm 0–3 người trên WiCount \canlam{kiểm tra Bảng V và cách chia train/test}. Các mô
hình này là mạng học sâu chạy trên máy tính; không có luật cảnh báo, bước xác nhận sau ngã hay đo nhịp thở.

\textbf{Gộp nhiều đường truyền.} Chang và cộng sự \cite{fusion2016} nhận dạng 7 hành động bằng CSI của card Intel 5300
với bốn cặp anten phát–thu. Gộp muộn (mỗi cặp một bộ phân loại, cộng xác suất) đạt 77,45\,\% và 86,96\,\% trên hai bộ
dữ liệu, cao hơn gộp sớm (ghép đặc trưng vào một bộ phân loại) với 70,24\,\% và 80,23\,\% (kiểm định chéo 10 lần).

\textbf{Đo nhịp thở bằng CSI của ESP32.} COVID‑Beat \cite{covidbeat} dùng một cặp ESP32, chọn các sóng mang con dao động
mạnh nhất, lọc dải rồi tìm đỉnh phổ; độ chính xác ước lượng nhịp thở từ 91\,\% đến 99\,\% trên 5 người. Burimas và cộng
sự \cite{appsci2024} đo nhịp thở khi ngủ bằng ESP32 ở các mức thở bình thường, nhanh và chậm \canlam{trích sai số công
bố}.

\textbf{Ngưỡng nhịp thở.} Thang NEWS2 \cite{news2} cho 3 điểm (mức nguy hiểm nhất của một chỉ số) khi nhịp thở $\le 8$
hoặc $\ge 25$ lần/phút, 0 điểm ở 12–20 lần/phút; đây là căn cứ chọn ngưỡng cảnh báo của hệ thống.

% =====================================================================================================================
\section{Hệ thống đề xuất}
\subsection{Phần cứng và thu dữ liệu}
Một ESP32‑C3 (TX) phát gói ESP‑NOW broadcast 100 gói/s trên một kênh cố định, băng thông 20\,MHz. Hai ESP32‑S3 (RX1,
RX2) đặt ở hai cạnh của khu theo dõi, cùng thu CSI của mỗi gói; CSI chỉ lấy phần LLTF, cho 52 sóng mang con hữu ích.
Mỗi RX lọc theo địa chỉ MAC của TX và đếm số khung của thiết bị khác trên kênh ($n_\text{filtered}$) để biết kênh có bị
chiếm hay không.

\subsection{Pipeline bốn tầng}
\begin{itemize}
  \item \textbf{Tầng 0} (window 6\,s, trượt 1\,s): phân loại phòng Trống / Hoạt động / Yên.
  \item \textbf{Tầng 1} (window 6\,s): 1 người hay $\ge 2$ người; phân vân thì chọn 1 để không tắt nhầm phát hiện ngã.
  \item \textbf{Tầng 2a} (window 3\,s, trượt 0,5\,s, chỉ khi 1 người): ngã, ngồi xuống, đứng dậy, nằm xuống, cúi nhặt,
    không có động tác. Bốn động tác thường được học riêng vì dễ nhầm với ngã.
  \item \textbf{Tầng 2b} (window 30\,s, trượt 5\,s, khi Yên): nhịp thở bằng lọc dải, PCA trên các sóng mang con, phổ
    Welch; chỉ chấp nhận đỉnh có độ nổi (đỉnh chia trung vị phổ) vượt ngưỡng.
\end{itemize}

\subsection{Đặc trưng và mô hình}
Thay vì 4 thống kê trên từng sóng mang con (212 cột mỗi board như \cite{natarajan2023}), window 6\,s dùng 19 cột: 4
thống kê trên RSSI, mức tín hiệu chung và hai thành phần chính PC1, PC2, cùng năng lượng dải thở và tỉ lệ gói hỏng.
Window 3\,s dùng 25 cột: năng lượng 5 dải tần 0,3–40\,Hz của PC1, đặc trưng hình dạng và số gói tốt. Mỗi tầng dùng
XGBoost 300 vòng, độ sâu 4 (6\,s) và 5 (động tác), trọng số cân bằng lớp.

\subsection{Gộp kết quả hai board}
Mỗi RX tự chạy mô hình trên CSI của nó và gửi xác suất về trạm chính. Xác suất ngã lấy \emph{max} (một board thấy là
đủ), xác suất nhiều người lấy \emph{min} (cả hai cùng thấy), mức hoạt động và động tác lấy trung bình, nhịp thở lấy
board có đỉnh phổ nổi nhất.

\subsection{Bộ cảnh báo theo sự kiện}
Ngã được báo khi thấy động tác ngã và 10\,s sau người chưa đứng dậy; huỷ khi thấy đứng dậy hoặc Hoạt động liên tục suốt
10\,s. Giây kênh bận ($n_\text{filtered} > 30$ hoặc $< 85$ gói/s) không được dùng để huỷ; hệ thống chờ kênh sạch tối đa
60\,s rồi vẫn báo. Bất động lâu: Yên liên tục $\ge 5$ phút. Thở bất thường: $> 25$ hoặc $< 8$ nhịp/phút. Ngưng thở: 3
cửa sổ liên tiếp không thấy nhịp sau khi đã đo được. Ngã mà không thấy cử động lẫn nhịp thở: cảnh báo khẩn cấp.

% =====================================================================================================================
\section{Thực nghiệm}
\subsection{Dữ liệu}
\begin{itemize}
  \item \textbf{Mô phỏng 3D}: đúng bố trí 1 TX + 2 RX; đi lại, ngồi, đứng, nằm, cúi nhặt, 3 kiểu ngã, ngã rồi giãy, thở
    8–36 nhịp/phút, nín thở, khách vào ra, nhiễu kênh; nhiễu và AGC hiệu chỉnh theo WiFall. 40 buổi, 4,3 giờ.
  \item \textbf{WiFall} \cite{wifall}: 928 đoạn 3\,s, dùng làm benchmark cho tầng động tác.
  \item \textbf{Dữ liệu thật của nhóm}: \canlam{mô tả phòng, số buổi, kịch bản, số lần lặp, cách lấy nhịp đáp án}.
\end{itemize}

\subsection{Cách đánh giá}
Chia train/test theo buổi (GroupKFold 5 lượt). Buổi test tính đặc trưng trượt từ dữ liệu thô, không dùng nhãn, như lúc
chạy thật. Ngã được tính là bắt được nếu có cảnh báo trong 30\,s.

\subsection{Baseline}
\canlam{Mỗi baseline chạy trên cùng dữ liệu, cùng cách chia:}
\begin{itemize}
  \item Phương pháp đầy đủ của \cite{natarajan2023}: 212 cột/board, bỏ phiếu RF + GB + XGBoost.
  \item Chỉ dùng RSSI; chỉ dùng 1 RX; gộp sớm đặc trưng hai board.
  \item Phát hiện ngã bằng ngưỡng năng lượng; nhịp thở bằng FFT trên một sóng mang con.
\end{itemize}

% =====================================================================================================================
\section{Kết quả}
\subsection{Bộ cảnh báo trên mô phỏng}
Bảng~\ref{tab:vong} cho kết quả qua ba vòng tối ưu (cách gộp muộn, có cờ kênh).

\begin{table}[t]
\caption{Kết quả theo sự kiện trên 40 buổi mô phỏng}
\label{tab:vong}
\centering\small
\begin{tabular}{@{}lccc@{}}
\toprule
 & Vòng 1 & Vòng 2 & Vòng 3 \\
\midrule
Ngã bắt được (/53) & 34 & 39 & \textbf{44} \\
Báo ngã giả (lần/giờ) & 14 & 6,7 & \textbf{5,1} \\
Bất động lâu bắt được (/8) & 3 & 5 & 5 \\
Thở bất thường bắt được (/21) & 8 & 2 & 6 \\
Nín thở bắt được (/12) & 5 & 1 & 1 \\
Lỗi đếm người nguy hiểm & 0,3\,\% & 0,4\,\% & 0,1\,\% \\
\bottomrule
\end{tabular}
\end{table}

\subsection{So với baseline}
\canlam{Bảng: độ chính xác / F1 / bắt ngã / báo giả của từng baseline và của nhóm, trung bình ± độ lệch qua nhiều seed;
kèm số cột và số nút cây.}

\subsection{Chạy trên ESP32‑S3}
Trên ESP32‑S3 N16R8 (240\,MHz), mô hình động tác cho cùng lớp với Python ở 40/40 window đối chiếu. Tính 25 cột mất
69\,ms, chạy mô hình 13\,ms mỗi window; với bước trượt 0,5\,s chiếm khoảng 16\,\% một nhân. \canlam{đo thêm khi vừa thu
CSI vừa tính, và hai mô hình 6\,s}

\subsection{Dữ liệu thật}
\canlam{Kết quả theo từng tình huống: đứng/ngồi/nằm yên ở giữa phòng, gần và xa đường sóng; ngã và động tác dễ nhầm;
nhịp thở so với metronome.}

% =====================================================================================================================
\section{Thảo luận và hạn chế}
\begin{itemize}
  \item Kết quả hiện chủ yếu trên mô phỏng; mô hình train trên mô phỏng chạy kém trên dữ liệu thật đầu tiên (hay gọi
    người ngồi yên là phòng trống), nên cần train lại trên dữ liệu thật.
  \item Nhịp thở và ngưng thở còn yếu; cần khoảng 40\,s mới báo ngưng thở, dài hơn định nghĩa y khoa ($\ge 10$\,s).
  \item Chỉ một phòng, một bố trí board; chưa đánh giá khi đổi phòng.
\end{itemize}

\section{Kết luận}
\canlam{Viết sau khi có kết quả baseline và dữ liệu thật.}

% =====================================================================================================================
\begin{thebibliography}{9}
\bibitem{natarajan2023} A. Natarajan, V. Krishnasamy, and M. Singh, ``A machine learning approach to passive human
  motion detection using WiFi measurements from commodity IoT devices,'' \emph{IEEE Trans. Instrum. Meas.}, vol. 72,
  Art. no. 2512510, 2023, doi: 10.1109/TIM.2023.3272374.
\bibitem{wifall} Z. Zhao, Z. Cai, T. Chen, X. Li, H. Li, Q. Chen, and G. Zhu, ``KNN‑MMD: Cross domain wireless sensing
  via local distribution alignment,'' arXiv:2412.04783, 2024.
\bibitem{csibert2} Z. Zhao, F. Meng, Z. Lyu, H. Li, X. Li, and G. Zhu, ``CSI‑BERT2: A BERT‑inspired framework for
  efficient CSI prediction and classification in wireless communication and sensing,'' \emph{IEEE Trans. Mobile
  Comput.}, 2025, doi: 10.1109/TMC.2025.3640420.
\bibitem{fusion2016} J.-Y. Chang, K.-Y. Lee, Y.-L. Wei, K. C.-J. Lin, and W. Hsu, ``We can `see' you via Wi‑Fi – WiFi
  action recognition via vision‑based methods,'' arXiv:1608.05461, 2016.
\bibitem{covidbeat} M. Atif, S. Muralidharan, H. Ko, and B. Yoo, ``COVID‑Beat: a low‑cost breath monitoring approach
  for people in quarantine during the pandemic,'' \emph{J. Comput. Des. Eng.}, vol. 9, no. 3, pp. 992–1006, 2022,
  doi: 10.1093/jcde/qwac037.
\bibitem{appsci2024} R. Burimas, T. Horanont, A. Thapa, and B. R. Lamichhane, ``Monitoring the sleep respiratory rate
  with low‑cost microcontroller Wi‑Fi in a controlled environment,'' \emph{Appl. Sci.}, vol. 14, no. 15, Art. no. 6458,
  2024, doi: 10.3390/app14156458.
\bibitem{news2} Royal College of Physicians, \emph{National Early Warning Score (NEWS) 2: Standardising the assessment
  of acute‑illness severity in the NHS}, London, 2017.
\bibitem{espcsi} Espressif Systems, ``esp‑csi,'' \url{https://github.com/espressif/esp-csi}.
\end{thebibliography}

\end{document}
